\section{UVM vs.\ VMM}
\label{sec:uvm-vs-vmm}

CUDA provides two distinct mechanisms for managing GPU memory beyond
\texttt{cudaMalloc}: Unified Memory (UVM, \texttt{cudaMallocManaged}) and the
Virtual Memory Management API (VMM,
\texttt{cuMemCreate}/\texttt{cuMemAddressReserve}/\texttt{cuMemMap}/\texttt{cuMemSetAccess}).
These mechanisms address different problems.

\begin{itemize}
  \item Unified Memory:
    \url{https://docs.nvidia.com/cuda/cuda-programming-guide/04-special-topics/unified-memory.html}
  \item Virtual Memory Management:
    \url{https://docs.nvidia.com/cuda/cuda-programming-guide/04-special-topics/virtual-memory-management.html}
\end{itemize}

\subsection{Why NVIDIA introduced each}

UVM addresses a programmer-burden problem. Prior to its introduction, every
host/device pointer pair required explicit \texttt{cudaMemcpy} calls to move
data between host and device code. \texttt{cudaMallocManaged} returns a
single pointer valid on both host and device, and the driver migrates pages
on demand as they are accessed. \texttt{cudaMemAdvise} and
\texttt{cudaMemPrefetchAsync} exist to improve performance beyond this
default behavior by supplying hints to the driver.

VMM addresses a different problem: it gives an application that already
knows its memory placement requirements a way to express them explicitly.
The CUDA Programming Guide motivates the API through a specific limitation
of the classic model: \texttt{cudaEnablePeerAccess} maps all of a process's
past and future allocations to a peer device, so sharing even a small number
of buffers incurs mapping overhead on every allocation the process ever
makes. VMM separates \texttt{cudaMalloc} into two steps.
\texttt{cuMemAddressReserve} reserves a virtual address range with no
physical backing; \texttt{cuMemCreate}, \texttt{cuMemMap}, and
\texttt{cuMemSetAccess} then bind physical memory to that range under
explicit, per-mapping access control.

\subsection{Representative workloads}

UVM is well suited to workloads with irregular, data-dependent access
patterns, where a hand-written copy schedule is impractical: graph
algorithms such as breadth-first search or PageRank on graphs too large or
too irregular to predict which region will be needed next, out-of-core
algorithms whose working set exceeds device memory, and early-stage kernel
development where correctness matters more than placement.

VMM is well suited to workloads that already know their placement
requirements and are penalized by not being able to express them: custom
allocators requiring in-place growth without copying data, key-value cache
managers in large language model inference, and multi-GPU communication
libraries such as NCCL and NVSHMEM, which share specific buffers across
devices rather than mapping entire address spaces.

\begin{center}
\begin{tabular}{@{}p{0.22\linewidth}p{0.37\linewidth}p{0.37\linewidth}@{}}
\toprule
 & \textbf{UVM} & \textbf{VMM} \\
\midrule
Core primitive & One pointer, driver migrates pages on fault & Reserve (VA), create (physical), map, and set-access, all explicit \\
Residency control & Advisory hints only & Deterministic: resident exactly where last mapped \\
Growth & Fixed size at allocation & Chunked, non-contiguous physical backing behind one contiguous virtual range \\
Cross-process sharing & Not a first-class primitive & First-class: handles are exportable and importable \\
\bottomrule
\end{tabular}
\end{center}

\subsection{Capabilities VMM provides that UVM does not}

\begin{itemize}
  \item \textbf{Non-contiguous physical backing behind one contiguous
    virtual address range.} A UVM allocation has a single, fixed physical
    footprint determined at call time. VMM allows separately allocated
    chunks to be mapped into one contiguous virtual range, which is the
    mechanism PyTorch's \texttt{expandable\_segments} allocator relies on.
  \item \textbf{Deterministic residency rather than a hint.}
    \texttt{cudaMemAdvise} is advisory and the driver may disregard it.
    \texttt{cuMemMap} and \texttt{cuMemUnmap} are not requests: once
    \texttt{cuMemMap} returns, the mapping exists. This property is what the
    \texttt{cudaremap} primitive in this project depends on, since it
    requires the caller to know precisely when a backing swap between
    device and host memory has completed.
  \item \textbf{Live, cross-process re-sharing of a single physical
    allocation.} A VMM handle can be mapped into more than one process's
    address space, and that mapping can change while other processes hold
    their own. UVM has no equivalent mechanism.
  \item \textbf{No page-fault safety net.} There is no fault handler for
    VMM-mapped ranges: a kernel accessing a virtual address that has been
    unmapped via \texttt{cuMemUnmap} terminates the process rather than
    faulting and recovering. UVM's entire mechanism is, by contrast, built
    around a fault handler. Any library built on VMM must therefore
    guarantee correctness of in-flight remap operations itself.
\end{itemize}

\subsection{Experiments to isolate the difference}

Establishing whether VMM's additional control is worthwhile on this
hardware, rather than UVM's driver heuristics already being sufficient,
requires direct measurement. The first experiment is as follows.

\textbf{Interrupt-driven paging versus advised prefetch with remap.} A
working set larger than device memory will be driven two ways: first with
UVM and no hints, relying entirely on fault-driven migration; second with
VMM, issuing \texttt{cudaremap} ahead of the kernel that requires the data.
Per-kernel stall time and total wall-clock time will be measured for each.
